\section{Almost bounded values of the Collatz map}

\begin{definition}[Collatz map]
\label{def:collatz-map}
\lean{CollatzMapAlmostBoundedValues.collatz}
\leanok
Define $T(n)=n/2$ when $n$ is even and $T(n)=3n+1$ when $n$ is odd.
\end{definition}

\begin{definition}[Attaining a value below a bound]
\label{def:collatz-attains-below}
\lean{CollatzMapAlmostBoundedValues.AttainsBelow}
\leanok
\uses{def:collatz-map}
For $f:\mathbb N\to\mathbb N$, say that the orbit from $n$ attains below $f(n)$ if some iterate $T^k(n)$ is strictly smaller than $f(n)$.
\end{definition}

\begin{definition}[Logarithmic density zero]
\label{def:collatz-log-density-zero}
\lean{CollatzMapAlmostBoundedValues.logWeightSum, CollatzMapAlmostBoundedValues.LogDensityZero}
\leanok
A set has logarithmic density zero when its normalized logarithmic harmonic weight tends to zero.
\end{definition}

\begin{theorem}[Almost all Collatz orbits attain almost bounded values]
\label{thm:collatz-almost-bounded}
\lean{CollatzMapAlmostBoundedValues.MainTheorem}
\notready
\uses{def:collatz-attains-below, def:collatz-log-density-zero}
For every $f:\mathbb N\to\mathbb N$ with $f(n)\to\infty$, the set of $n$ whose Collatz orbit never falls below $f(n)$ has logarithmic density zero.
\end{theorem}

\begin{proof}
The definitions and statement are formalized; the proof of the main theorem is currently a \texttt{sorry}.
\end{proof}
